\section{Aprendizado Baseado em Simulação (SBL)}
\label{sec:sbl}

O Aprendizado Baseado em Simulação (\textit{Simulation-Based Learning} — SBL) no escopo da odontologia não se restringe a uma ferramenta acessória, mas configura um paradigma epistemológico robusto ancorado nas ciências cognitivas. Sustentado pelos fundamentos da \textit{teoria da aprendizagem experiencial} de David Kolb, o SBL postula que o conhecimento clínico e procedimental é ativamente construído por meio da dialética entre apreensão (experiência concreta e conceitualização abstrata) e transformação (observação reflexiva e experimentação ativa). 

A literatura acadêmica recente e de alto impacto reitera, com base em ensaios controlados randomizados de larga escala, que a imersão mediada por gêmeos digitais transpõe limitações inerentes ao currículo clássico. Como investigado por Robles et al. (2023) no contexto da validação de simulações abertas, os estudantes expostos a ciclos de SBL apresentaram incrementos significativos tanto na coordenação visomotora fina quanto na resiliência cognitiva em cenários clínicos não programados~\cite{robles2023opentwins}. A introdução sistemática da simulação ciber-física altera a dinâmica da \textit{zona de desenvolvimento proximal} de Vygotsky, com o sistema atuando como um tutor virtual que nivela o desafio à proficiência dinâmica do discente.

\destaque{A virtude primária do SBL com gêmeos digitais reside na "prática deliberada" (Anders Ericsson). A capacidade de reinicializar indefinidamente uma patologia virtual complexa permite ao aluno atingir a sobreaprendizagem (\textit{overlearning}), essencial para a automação de habilidades psicomotoras finas sob estresse clínico.}

Uma meta-análise abrangente delineada por Shen et al. (2024) documentou que estudantes que completaram rotinas de treinamento baseadas em RV em módulos de exodontia e periodontia relataram índices de ansiedade pré-clínica substancialmente inferiores, além de uma taxa de preservação de tecido sadio até 30\% superior durante suas primeiras intervenções em pacientes reais, corroborando a eficácia tangível do SBL na translação de habilidades pré-clínicas~\cite{shen2024effectiveness}.

Ademais, o aprendizado baseado em gêmeos digitais converte os erros de um passivo ético e financeiro em um ativo pedagógico. Em odontologia clássica, o erro iatrogênico em paciente é um evento a ser evitado a todo custo, o que paradoxalmente subtrai do aluno a oportunidade valiosa de aprender o manejo de complicações. No ambiente hiper-realista da simulação, a perfuração de uma raiz ou o sobreaquecimento ósseo desencadeiam protocolos virtuais de resgate, instruindo o aluno sobre fluxos de contingência em tempo real sem comprometer a integridade de seres humanos~\cite{bdj2026predictive}.
